\section{Redundancies six and seven: complete and gated results}
\label{sec:r67}

This appendix takes the marked polar-escape theorem as input and discharges
its first two pointed lower packages.  The bottleneck is exhaustiveness: the
persistent rank-two carrier, modular nucleus, self-collision, and transverse
cases must account
for every polar line.  The resulting structural arguments, together with the
bounded certificates, prove Theorems~\ref{thm:r6} and~\ref{thm:r7}.

The numerical mechanism is compact:
\[
\begin{array}{c|c|c|c|c}
\text{level}&\text{stages}&\text{marker-scheme degrees}
 &\text{terminal bad points}&\text{first field gate}\\ \hline
\mathrm{R6}&1&13&18&29\\
\mathrm{R7}&2&12,14&24&37.
\end{array}
\]
The corresponding contained residues are the persistent carrier together with
the binary nucleus line at R6 and its central-point lift at R7.

\subsection{Redundancy six}

For
\[
 f=(a_0:\cdots:a_5)\in\PP(\Gamma^5E),\qquad
 W_f=\ker
 \begin{pmatrix}
 a_0&a_1&a_2&a_3&a_4\\
 a_1&a_2&a_3&a_4&a_5
 \end{pmatrix},
\]
the system $W_f$ is a net of quartics off the quintic normal rational
curve.  It is split-free exactly when it has no completely split
squarefree quartic.

\begin{proposition}[Persistent net and orbit law]\label{prop:r6-persistent}%
\coverage{fragment}%
\lean{PRSRedundancySixSeven.PersistentModularFamilyData.classified_card_doubled,PRSPolarInduction.fifthPower_sigmaInversionOrbitCount}%
If $\gcd(W_f)=Q$ has degree two, then
\[
       W_f=Q\Sym^2E^\vee.
\]
The split-free persistent locus has size $q(q+1)^2/2$.  On the tangent fiber
the unipotent coordinate is $x\mapsto x+5u$; on the irreducible
quadratic fiber the norm-one torus acts through $z\mapsto z^5$, with
inversion and coefficient Frobenius acting on $T/T^5$.
\end{proposition}

\begin{proof}
For fixed $Q$, the syndromes annihilating
$Q\Sym^3E^\vee$ form a projective line.  If $Q$ is irreducible, all
$q+1$ points remain; if $Q=\lambda^2$, one endpoint has dependent
Hankel rows, leaving $q$ points.  Hence
\[
 q(q+1)+\frac{q(q-1)}2(q+1)=\frac{q(q+1)^2}{2}.
\]
Normalize a square to $U^2$.  The nondegenerate fiber is
$[e_4+xe_5]$, and direct degree-five substitution gives
$x\mapsto x+5u$.  For an irreducible $Q$, its stabilizer identifies
the $q+1$-point fiber with
$T=\{z\in\F_{q^2}^\times:z^{q+1}=1\}$ and acts by fifth powers; the
normalizer inverts $z$, while coefficient Frobenius multiplies
$T/T^5$ by the characteristic.  This proves the stated orbit law and
the characteristic-five tangent splitting.
\end{proof}

\begin{proposition}[First-polar line and secant degree]
\label{prop:r6-secant}%
\coverage{fragment}%
\lean{PRSPolarInduction.sequenceContraction_agrees_with_finite}%
For $r\in\F_q$ put
\[
 b(r)=(a_1-ra_0,\ldots,a_5-ra_4),\qquad
 b(\infty)=(a_0,\ldots,a_4).
\]
Then
\[
 (T-rU)g\in W_f\quad\Longleftrightarrow\quad g\in W_{b(r)}.
\]
If $W_f$ has trivial gcd, its polar line meets the lower secant cubic
in degree at most three.
\end{proposition}

\begin{proof}
The first equivalence follows by expanding the two Hankel equations of
$(T-rU)g$.  Let
\[
C_f=\begin{pmatrix}
a_0&a_1&a_2&a_3\\
a_1&a_2&a_3&a_4\\
a_2&a_3&a_4&a_5
\end{pmatrix}.
\]
The lower catalecticant equals
\[
C_f\begin{pmatrix}
-r&0&0\\1&-r&0\\0&1&-r\\0&0&1
\end{pmatrix}.
\]
Cauchy--Binet gives the secant determinant
\[
D_f(r)=-\Delta_{012}r^3+\Delta_{013}r^2
       -\Delta_{023}r+\Delta_{123}.
\]
It vanishes identically exactly when $\operatorname{rank}C_f\leq2$.  On the
rank-two open this is equivalent to $W_f$ having quadratic gcd; rank one is
the normal-rational-curve boundary.  In the trivial-gcd case it therefore has
at most three zeros.
\end{proof}

\begin{proposition}[Complete lower trivial-gcd carrier]
\label{prop:r6-lower-carrier}%
\coverage{absent}%
Let \(c\in\PP(\Gamma^4E)\) and suppose that the cubic pencil \(W_c\)
has trivial gcd.  One of the following applies over an algebraic
closure, and the last three alternatives exhaust the cases in which
the \(S_3\) package in (i) is unavailable:
\begin{enumerate}[label=\textup{(\roman*)}]
\item the cubic map is separable with geometric monodromy \(S_3\), and
      its off-diagonal fiber square is geometrically integral;
\item the characteristic is neither two nor three and the pencil is cyclic, in
      which case \(c\) lies on the closure of
      \[
      \left\{\operatorname{diag}(1,4,6,4,1)^{-1}[Q^2]:
      [Q]\in\PP(\Sym^2E^\vee)\right\};
      \]
\item the characteristic is two and \(c\) lies on
      \(\Pi_{\mathrm{cyc}}=\PP\langle e_1,e_2,e_3\rangle\);
\item the characteristic is three and \(c\) lies on
      \(\operatorname{Join}(e_2,C_4)\), the closure of the wild cyclic
      pencils together with the inseparable nucleus point.
\end{enumerate}
Thus these displayed tame, binary, and ternary loci exhaust the
trivial-gcd lower pencils for which the \(S_3\) package is unavailable.
\end{proposition}

\begin{proof}
Proposition~\ref{prop:r5-incidence} associates to \(W_c\) a degree-three
map.  If it is separable, its transitive geometric monodromy is
\(C_3\) or \(S_3\).  In the \(S_3\) case two-transitivity makes the
off-diagonal fiber square geometrically integral, giving (i).
In the cyclic case, Lemma~\ref{lem:cyclic} identifies the two
ramification points with a binary quadratic \(Q\).  Substitution in the
two Hankel equations gives the divided-power syndrome
\(\operatorname{diag}(1,4,6,4,1)^{-1}[Q^2]\) whenever the displayed
diagonal is invertible.  This proves (ii) in characteristics other than
two and three.

The same coefficient calculation can be made without division.  In
characteristic two it is
\[
 g\cdot e_2=
 [0:\delta\gamma:\alpha\delta+\beta\gamma:\alpha\beta:0],
 \qquad
 g=\begin{pmatrix}\alpha&\beta\\ \gamma&\delta\end{pmatrix}.
\]
Its image is dense in the plane
\(\PP\langle e_1,e_2,e_3\rangle\), so its closure is exactly the
binary locus in (iii).  In characteristic three every separable
cyclic pencil is wild.  Lemma~\ref{lem:cyclic} puts it in the
\(\PGL_2\)-orbit of \(e_2-ae_4\).  Here \(e_2\) is fixed and the orbit
closure of \(e_4\) is the quartic normal rational curve \(C_4\);
hence the complete wild closure is
\(\operatorname{Join}(e_2,C_4)\).  The only inseparable trivial-gcd
pencil is the all-cube nucleus pencil, namely its vertex \(e_2\), by
Proposition~\ref{prop:r5-bridge}.  This proves (iv) and exhausts the
inseparable case.  The three alternatives are therefore exhaustive.
\end{proof}

\begin{proposition}[Cyclic, collision, and terminal bounds]\label{prop:r6-degrees}%
\coverage{absent}%
For a trivial-gcd net:
\begin{enumerate}[label=\textup{(\roman*)}]
\item in characteristic different from two and three, the lower tame cyclic
carrier is a binomially rescaled projected quadratic Veronese surface
of degree four and
contains no polar line;
\item the pointed ramification condition has degree at most six and
does not vanish identically;
\item on an $S_3$ lower pencil, at most eighteen rational points of the
off-diagonal fiber square fail to encode a completely split member avoiding a
specified marked root.
\end{enumerate}
\end{proposition}

\begin{proof}
In syndrome coordinates the tame cyclic surface is
\[
\left\{\operatorname{diag}(1,4,6,4,1)^{-1}[Q^2]:
[Q]\in\PP(\Sym^2E^\vee)\right\}.
\]
When $2$ and $3$ are invertible this is the projected quadratic
Veronese surface, of degree four, and contains no line.
Characteristic three has a different
wild closure, treated in Proposition~\ref{prop:r6-modular}.
For (ii), Lemma~\ref{lem:uniform-collision} with \(j=6\) gives the
degree-six bound except possibly in characteristic two.  In that
characteristic, inseparability would identify \(W_f\), after a change
of coordinates, with the full pure-square space
\(\langle T^4,T^2U^2,U^4\rangle\).  Its annihilator is
\(\langle e_1^\vee,e_3^\vee\rangle\).  For the two consecutive Hankel
rows this would require simultaneously
\[
 a_0=a_2=a_4=0,\qquad a_1=a_3=a_5=0,
\]
contrary to \(f\ne0\).  Thus the bound holds in every characteristic.

Finally, Proposition~\ref{prop:r5-count} and Lemma~\ref{lem:r5-branch} show
that at most twelve rational points encode nonsplit members.
Base-point-freeness gives a unique pencil member through the specified marker,
and it contributes at most six further points.  Thus the terminal bad set has
size at most \(18\), and a permitted split member exists as soon as
\[
       q+1-2\sqrt q>18.
\]
The first prime-power solution is $29$.
\end{proof}

\begin{proposition}[Modular contained components]\label{prop:r6-modular}%
\coverage{fragment}%
\lean{PRSPolarInduction.mem_modularContractionKernel_iff}%
In characteristic three the lower wild carrier is the degree-four cone
\[
       \operatorname{Join}(e_2,C_4),
\]
and no polar line of a trivial-gcd net is contained in it.  In
characteristic two the cyclic carrier is the plane
\[
       \Pi_{\mathrm{cyc}}=\PP\langle e_1,e_2,e_3\rangle,
\]
whose unique contained polar-line component is
\[
       \cN_3=\PP\langle e_2,e_3\rangle.
\]
\end{proposition}

\begin{proof}
Every line on the wild cone is a ruling.  By equivariance it suffices
to compare the two consecutive rows on one ruling; in the frame
$\langle e_2,e_4\rangle$ the conditions
\[
(a_0,\ldots,a_4),\ (a_1,\ldots,a_5)\in\langle e_2,e_4\rangle
\]
force the overlapping coordinates successively to vanish.  Undoing
the frame gives $f=\mu\nu(t)$, including $f=\mu e_5$ at infinity.
These are fixed-factor/rank-boundary syndromes, so no trivial-gcd
polar line is contained in the wild cone.

In characteristic two, direct substitution gives
\[
g\cdot e_2=
[0:\delta\gamma:\alpha\delta+\beta\gamma:\alpha\beta:0],
\]
whose closure is $\Pi_{\mathrm{cyc}}$.  The consecutive-row overlap
is now literal: the first row lies in the plane exactly when
$a_0=a_4=0$, and the second exactly when $a_1=a_5=0$.  Hence a polar
line lies in the plane exactly when
\[
a_0=a_1=a_4=a_5=0,
\]
namely when $f\in\PP\langle e_2,e_3\rangle$.  This proves both
existence and uniqueness of the contained binary component.
\end{proof}

\begin{proposition}[Binary nucleus arithmetic]\label{prop:r6-nucleus}%
\coverage{absent}%
The scheme $\cN_3(\F_q)$ is one projective-semilinear orbit of $q+1$
points.  It is split-free over $q=2^m$ exactly when $m$ is odd.
\end{proposition}

\begin{proof}
For the representative $e_3$,
\[
       W_{e_3}=\langle1,t,t^4\rangle.
\]
If \(C=0\), the member \(A+Bt+Ct^4\) has degree at most one and
cannot have four distinct roots.  Assume \(C\ne0\).  The root
differences of $A+Bt+Ct^4$ then lie in the kernel of
$u\mapsto u^4+(B/C)u$.  Four distinct rational roots exist exactly
when all roots of $u^3=B/C$ are rational, equivalently when
$3\mid q-1$, which for $q=2^m$ means $m$ is even.  The stabilizer of
$e_3$ is the Borel of order $q(q-1)$, so the orbit has $q+1$ points;
its lower-unipotent transforms and endpoint are precisely
$\cN_3(\F_q)$.
\end{proof}

\pagebreak[3]
\begin{proposition}[Exhaustive R6 contained-component classification]
\label{prop:r6-contained}%
\coverage{absent}%
For a redundancy-six syndrome, the alternatives in
Definition~\ref{def:contained} are exhausted as follows:
\begin{enumerate}[label=\textup{(\roman*)}]
\item a noninjective contraction map is the rank-one
      normal-rational-curve boundary and is shallow;
\item a polar line contained in the lower secant cubic comes from the
      upper rank-two catalecticant locus; after its shallow rational
      secants are removed, this is the persistent tangent and
      conjugate-secant locus;
\item no trivial-gcd polar line is contained in the tame cyclic
      surface or the characteristic-three wild cone;
\item in characteristic two, containment in the cyclic plane occurs
      exactly for \(f\in\cN_3=\PP\langle e_2,e_3\rangle\);
\item on the remaining trivial-gcd locus the collision divisor
      is not identically zero.
\end{enumerate}
Consequently the complete nonshallow contained list is the persistent
locus together with \(\cN_3\) in characteristic two.  There is no
unlisted contained component.
\end{proposition}

\begin{proof}
The noninjective and secant-cubic alternatives are
Proposition~\ref{prop:contained-rank-two} and
Proposition~\ref{prop:r6-secant}.  The tame cyclic and collision
assertions are Proposition~\ref{prop:r6-degrees}; the wild and binary
specializations are Proposition~\ref{prop:r6-modular}.
Proposition~\ref{prop:r6-lower-carrier} proves that these are the
complete trivial-gcd lower failure loci, rather than a selected list
of examples.  Together they are all the clauses of the terminal carrier in
Definition~\ref{def:contained}: contraction
indeterminacy, containment in each component of the lower bad scheme,
and identical collision.  Hence the list is exhaustive.
\end{proof}

\begin{proposition}[R6 one-step exhaustion]\label{prop:r6-one-step}%
\coverage{conditional_deduction}%
\lean{PRSRedundancySixSeven.redundancySixHighFieldSynthesis}%
Assume \(q\geq29\).
Then every split-free redundancy-six syndrome lies in the persistent
quadratic-gcd locus or, when \(q=2^m\) with odd \(m\), in
\(\cN_3(\F_q)\).
\end{proposition}

\begin{proof}
\emph{Contained nets.}
The complete contained branch is
Proposition~\ref{prop:r6-contained}.
Lemma~\ref{lem:basepointfree-off-persistent} puts every common-factor system
in the persistent rank-two locus or on its shallow rank-one boundary.  Among
the rank-two points, removing the shallow rational secants leaves exactly the
tangent and conjugate-secant families.  It remains to treat a trivial-gcd net.

\emph{Transverse exclusions.}
Its first-polar line is not contained in the lower secant cubic, by
Proposition~\ref{prop:r6-secant}.  The pullback of that cubic has degree
at most three.  The tame cyclic surface has degree four and contains
no polar line.  In characteristic three its wild replacement has the
same degree and no contained trivial-gcd polar line; in characteristic
two the replacement is a plane, whose only contained component is
\(\cN_3\), by Proposition~\ref{prop:r6-modular}.  Finally,
Proposition~\ref{prop:r6-degrees} gives a nonzero collision
divisor of degree at most six.  Outside the contained binary component,
the excluded parameter scheme therefore has degree
\[
                         d\leq3+4+6=13.
\]

\emph{Bottom pencils.}
For every remaining marker, the lower syndrome lies outside \(\cP_5\) and the
cyclic locus.  Lemma~\ref{lem:basepointfree-off-persistent}
makes its cubic pencil base-point-free, and Lemma~\ref{lem:s3} supplies the
geometrically integral off-diagonal curve.  If no split member avoids the
marker, Proposition~\ref{prop:r6-degrees}(iii) bounds its rational points by
eighteen.  Aubry--Perret therefore supplies an allowed member when
\[
 q+1-2\sqrt q>18,
\]
whose first prime-power solution is \(29\).
Together with the one-step parameter package above, of degree \(d=13\), this
terminal bad-point bound meets Theorem~\ref{thm:induction}, which produces a
split squarefree quartic.

\emph{The modular residue.}
This is one field-order threshold, not three characteristic-dependent
estimates: the first characteristic-three and characteristic-two
prime powers above it happen to be \(81\) and \(32\), respectively.
On \(\cN_3\),
Proposition~\ref{prop:r6-nucleus} gives the asserted odd-degree
split-free orbit and proves shallowness in even extension degree.
\end{proof}

The covering-radius input is complete.  Its exact high-rate endpoint is
\(q=8\), since \(6=\lfloor8/2\rfloor+2\).  Thus Seroussi--Roth
nonextendability and Dür's equivalence give
\(\rho(\PRS(q-5))=5\) for \(q\geq8\), while Certificate R6 checks it
directly at \(q=7\) and independently confirms \(q=8\).

\begin{theorem}[Redundancy six]\label{thm:r6}%
\coverage{conditional_deduction}%
\lean{PRSRedundancySixSeven.redundancySixAllFieldSynthesis,PRSRedundancySixSevenCertificate.redundancySix_count_exhaustion,PRSUniformCoveringRadius.SeroussiRothDuerRadiusInput.seroussiRothDimensionRange_six_eight,radiusRange_six_eight_of_externalSeroussiRothDuer}%
\uses{thm:r5}%
\imports{duer,seroussi-roth}%
\evidence{certificate-r6}%
\proves{redundancy-six-classification}%
For every prime power $q\geq7$, the deep syndromes of $\PRS(q-5)$ are
the persistent locus of size $q(q+1)^2/2$, together with:
\begin{enumerate}[label=\textup{(\roman*)}]
\item respectively $18,11,4,2,1$ exceptional $\PGL_2(q)$ orbits at
      $q=7,8,9,11,13$;
\item one characteristic-two nucleus orbit of size $q+1$ when
      $q=2^m$ and odd $m\geq5$.
\end{enumerate}
The lower bound \(m\geq5\) avoids double counting: the remaining odd
case \(m=3\), namely \(q=8\), is already the
\(|\mathcal O|=9\) row of the exceptional table.
There are no other exceptional orbits.  The complete persistent orbit
law is $T/T^5$ modulo inversion and Frobenius, with a tangent
fixed/nonzero split exactly in characteristic five.
At \(q=7,8,9,11,13\), the exceptional direction totals are
\[
                  5152,\quad4713,\quad1800,\quad792,\quad546,
\]
and the corresponding total numbers of deep syndrome directions are
\[
                  5376,\quad5037,\quad2250,\quad1584,\quad1820.
\]
\end{theorem}

\begin{proof}
Propositions~\ref{prop:r6-contained} and~\ref{prop:r6-one-step} give
the complete conceptual exhaustion for every prime power \(q\geq29\).
Certificate R6 exhausts
\[
q\in\{7,8,9,11,13,16,17,19,23,25,27\},
\]
finds exceptional projective-orbit counts $18,11,4,2,1$ only at the
five displayed fields, and no further orbit at the remaining fields.
Together with the separate covering-radius argument, this covers every prime power
$q\geq7$.
\end{proof}

The listed sources are disjoint.  The persistent locus has quadratic
gcd, every finite exceptional row has trivial gcd, and the separate
nucleus term occurs only for binary fields beyond the finite table.

The five small tables are exhaustive computations.  For an exceptional
net \(W_f\), let \(I_f(r)\) be the number of irreducible cubic members
of the quotient pencil \(W_{b(r)}\), including the infinity chart, and
define its shared-root pair count by
\[
                 E_f=\sum_{r\in\PP^1(\F_q)}\binom{I_f(r)}2.
\]
In odd characteristic let
\[
                 F_f(T,U)=\sum_{i=0}^5a_iT^{5-i}U^i
\]
be the binary quintic attached to \(f=(a_0:\cdots:a_5)\), and let
\(\tau_5(f)\) be its rational factorization type; for example,
\(1^3+2\) denotes a triple rational factor and an irreducible
quadratic factor.  Their semilinear orbit counts are
\(18,5,2,2,1\); \(E_f\), together in odd characteristic with
\(\tau_5(f)\), separates the classes up to precisely Frobenius fusion.
Table~\ref{tab:r6-exceptional} prints
the canonical representatives, stabilizers, orbit sizes, and fusion
lengths.

% Generated by supplement/build_r6_paper_table.py from the
% hash-pinned supplement/CLASSIFICATION-RECORDS.json.
\begin{longtable}{c l r r c r l}
\caption{Exceptional redundancy-six classes.  Each row is one
projective-semilinear class.  The columns give a canonical syndrome
representative, the size and stabilizer order of each constituent
\(\PGL_2(q)\)-orbit, the number \(\phi\) of such orbits fused by
coefficient Frobenius, the shared-root pair count \(E\), and the binary-quintic
factor type \(\tau_5\) in odd characteristic.}
\label{tab:r6-exceptional}\\
\toprule
\(q\) & representative & \(|\mathcal O|\) & \(|G_f|\) & \(\phi\) & \(E\) & \(\tau_5\)\\
\midrule
\endfirsthead
\toprule
\(q\) & representative & \(|\mathcal O|\) & \(|G_f|\) & \(\phi\) & \(E\) & \(\tau_5\)\\
\midrule
\endhead
\bottomrule
\endfoot
7 & \([0:0:0:1:0:1]\) & 168 & 2 & 1 & 7 & \(1^3+2\)\\
7 & \([0:0:1:0:3:1]\) & 336 & 1 & 1 & 19 & \(1^2+1+2\)\\
7 & \([0:0:1:0:3:2]\) & 336 & 1 & 1 & 15 & \(1^2+3\)\\
7 & \([0:1:0:0:0:1]\) & 168 & 2 & 1 & 10 & \(1+2+2\)\\
7 & \([0:1:0:0:2:0]\) & 112 & 3 & 1 & 33 & \(1+1+3\)\\
7 & \([0:1:0:1:1:2]\) & 336 & 1 & 1 & 9 & \(1+4\)\\
7 & \([0:1:0:1:1:5]\) & 336 & 1 & 1 & 9 & \(1+1+3\)\\
7 & \([0:1:0:1:3:0]\) & 336 & 1 & 1 & 21 & \(1+1+1+2\)\\
7 & \([0:1:0:1:3:3]\) & 336 & 1 & 1 & 7 & \(1+1+3\)\\
7 & \([0:1:0:1:3:6]\) & 336 & 1 & 1 & 14 & \(1+1+3\)\\
7 & \([0:1:0:3:0:5]\) & 168 & 2 & 1 & 11 & \(1+4\)\\
7 & \([0:1:0:3:0:6]\) & 168 & 2 & 1 & 20 & \(1+4\)\\
7 & \([0:1:0:3:2:1]\) & 336 & 1 & 1 & 13 & \(1+1+3\)\\
7 & \([1:0:0:0:1:2]\) & 336 & 1 & 1 & 7 & \(1+4\)\\
7 & \([1:0:0:0:1:3]\) & 336 & 1 & 1 & 23 & \(5\)\\
7 & \([1:0:0:3:1:3]\) & 336 & 1 & 1 & 13 & \(1+4\)\\
7 & \([1:0:1:0:5:2]\) & 336 & 1 & 1 & 8 & \(1+4\)\\
7 & \([1:0:1:0:6:2]\) & 336 & 1 & 1 & 8 & \(5\)\\
8 & \([0:0:0:1:0:0]\) & 9 & 56 & 1 & 0 & --\\
8 & \([0:1:1:0:2:1]\) & 504 & 1 & 3 & 20 & --\\
8 & \([0:1:1:0:2:5]\) & 504 & 1 & 3 & 19 & --\\
8 & \([1:0:0:1:3:6]\) & 504 & 1 & 3 & 14 & --\\
8 & \([1:0:1:3:5:2]\) & 168 & 3 & 1 & 54 & --\\
9 & \([0:1:0:0:0:4]\) & 180 & 4 & 2 & 28 & \(1+4\)\\
9 & \([0:1:0:1:4:2]\) & 720 & 1 & 2 & 29 & \(1+4\)\\
11 & \([0:0:0:1:0:0]\) & 132 & 10 & 1 & 60 & \(1^3+1^2\)\\
11 & \([0:1:0:2:0:10]\) & 660 & 2 & 1 & 30 & \(1+4\)\\
13 & \([0:1:0:0:0:2]\) & 546 & 4 & 1 & 60 & \(1+4\)\\
\end{longtable}
\noindent For \(q=8\), integers are binary polynomial-basis codes
with \(\alpha^3=\alpha+1\); for \(q=9\), they are ternary
polynomial-basis codes with \(\alpha^2=-1\).  Thus
\(\sum\phi=(18, 11, 4, 2, 1)\) and
\(\sum\phi|\mathcal O|=(5152, 4713, 1800, 792, 546)\),
in the field order \(7,8,9,11,13\).


\subsection{Redundancy seven}

For
\[
f=(a_0:\ldots:a_6)\in\PP(\Gamma^6E),\qquad
W_f=\ker
\begin{pmatrix}
a_0&a_1&a_2&a_3&a_4&a_5\\
a_1&a_2&a_3&a_4&a_5&a_6
\end{pmatrix},
\]
the system $W_f$ is a four-dimensional web of quintics.  For an affine
marker $r$ the contractions are
\[
 b(r)=(a_1-ra_0,\ldots,a_6-ra_5).
\]
Coefficient collection gives
\[
(t-r)g\in W_f\quad\Longleftrightarrow\quad g\in W_{b(r)};
\]
the lift is squarefree exactly when \(g\) is squarefree and \(g(r)\ne0\).

On the rank-two Hankel-row open, the contracted net \(W_{b(r)}\) has quadratic
gcd exactly when its three-row catalecticant has rank at most two.  The
rank-one boundary is part of the persistent contained case.  Unless the
first-polar line is contained in that rank-two locus, the remaining parameters
form the degree-at-most-three intersection of
Proposition~\ref{prop:r6-secant}; the contained case is
Corollary~\ref{cor:r7-contained}.

\begin{proposition}[Complete two-marker lower package]
\label{prop:r7-pointed}%
\coverage{conditional_deduction}%
\lean{PRSPolarInduction.LowerCoverStratum,CoherentPolarInput.splitFree_implies_persistent_or_modular}%
Let \(W_h\) be a trivial-gcd quartic net and let \(x\) be a prescribed
forbidden root.  Assume in characteristic two that \(h\notin\cN_3\).
For every prime power \(q\geq37\), the net contains a split squarefree
quartic avoiding \(x\).
\end{proposition}

\begin{proof}
\emph{Choosing the second marker.}
By equivariance choose coordinates in which \(x\) is finite; all
parameter equations below are then homogenized in \(s\), so the
infinity chart is included.
Choose the second contraction marker \(s\).  The unavailable values of
\(s\) comprise: \(s=x\), of degree one; at most three intersections
with the lower secant cubic, by
Proposition~\ref{prop:r6-secant}; at most four intersections with the
cyclic surface or its characteristic-three wild replacement; and at
most six points of the collision divisor, by
Propositions~\ref{prop:r6-degrees} and~\ref{prop:r6-modular}.  In
characteristic two the cyclic
replacement is a plane, so its pullback has degree at most one unless
the polar line is contained in it, the excluded case
\(h\in\cN_3\), by Proposition~\ref{prop:r6-modular}.  A lower pencil with
any fixed factor already lies in the degree-three persistent intersection, by
Lemma~\ref{lem:basepointfree-off-persistent}.  Hence
the complete second-step parameter
scheme has degree at most
\[
                         1+3+4+6=14<q+1.
\]
Choose \(s\) outside it.

\emph{The \(S_3\) branch.}
The bottom syndrome lies outside the persistent carrier, so
Lemma~\ref{lem:basepointfree-off-persistent} makes its pencil base-point-free;
outside the cyclic locus its cover has geometric monodromy \(S_3\).  Its
off-diagonal curve is geometrically integral of arithmetic genus one.
Proposition~\ref{prop:r5-count} and Lemma~\ref{lem:r5-branch} show that at most
twelve rational points encode nonsplit members.  Base-point-freeness gives at
most one pencil member through each of the distinct markers \(x,s\), and these
contribute at most twelve further points.  Thus at most \(24\) rational points
of the bottom off-diagonal curve fail to encode a split member avoiding both
markers.  Aubry--Perret supplies an allowed member when
\[
 q+1-2\sqrt q>24,
\]
whose first prime-power solution is \(37\).  The cubic
lifts through \(s\) to a split squarefree quartic in \(W_h\) avoiding
\(x\).
\end{proof}

\begin{proposition}[Collision divisor]\label{prop:r7-collision}%
\coverage{absent}%
After fixed factors are removed, the collision divisor of the
moving $g^3_5$ has degree at most eight.
\end{proposition}

\begin{proof}
This is Lemma~\ref{lem:uniform-collision} with \(j=7\):
the series is separable and \(2j-6=8\).
\end{proof}

The characteristic-two nucleus line at redundancy six has exactly one
coherent lift, the central point $e_3$, with
\[
W_{e_3}=\langle1,t,t^4,t^5\rangle.
\]
It is split-free for odd \(m\), and hence deep whenever the separate
radius-six gate holds.  For even \(m\), the homogenization of
\(t^4+t\) has the five roots \(\PP^1(\F_4)\).

\begin{proposition}[Central binary lift]\label{prop:r7-central}%
\coverage{absent}%
The complete inverse image of the R6 nucleus line under consecutive
sextic contraction is $[e_3]$.  It is split-free over $\F_{2^m}$
exactly when $m$ is odd.
\end{proposition}

\begin{proof}
Putting both consecutive rows in
$\PP\langle e_2,e_3\rangle$ gives
$a_0=a_1=a_2=a_4=a_5=a_6=0$.  Thus only $e_3$ remains, with
$W_{e_3}=\langle1,t,t^4,t^5\rangle$.  For odd $m$, contraction at a
root of a hypothetical split quintic would give a split quartic on the
R6 nucleus line, contradicting Proposition~\ref{prop:r6-nucleus}.  For
even $m$, $\F_4\subset\F_q$ and the homogenization of $t^4+t$ has its
four affine $\F_4$ roots and the simple root at infinity.
\end{proof}

\begin{corollary}[Exhaustive R7 contained-component classification]
\label{cor:r7-contained}%
\coverage{absent}%
The one-step contained-component assertion from redundancy seven to
redundancy six is exhaustive in every characteristic.  A noninjective
contraction is the shallow rank-one
boundary.  Containment in the lower persistent locus gives upper
catalecticant rank at most two; after the shallow rational secants are
removed, these are exactly the persistent tangent and conjugate-secant
families.  The complete lift of the only additional lower component,
the characteristic-two nucleus line, is the central point \(e_3\).
The remaining collision divisor is nonzero.  Thus the
nonshallow contained list is precisely the persistent locus and, in
characteristic two, \(e_3\); there is no further contained component.
\end{corollary}

\begin{proof}
Apply Proposition~\ref{prop:contained-rank-two} with $n=6$.  The three
rational types of the quadratic recurrence are split squarefree,
repeated rational, and irreducible.  They give respectively the shallow
rational secant, tangent, and conjugate-secant loci.  The remaining
lower binary nucleus line has the separate complete lift computed in
Proposition~\ref{prop:r7-central}.  Proposition~\ref{prop:r7-collision}
shows that the collision divisor is nonzero.  These are respectively
the contraction-indeterminacy, lower-component-containment, and
identical-collision alternatives of Definition~\ref{def:contained},
so the list is exhaustive.
\end{proof}

\pagebreak[3]
\begin{theorem}[Redundancy-seven classification]
\label{thm:r7}%
\coverage{conditional_deduction}%
\lean{PRSRedundancySixSeven.redundancySevenAllFieldSynthesis,PRSRedundancySixSevenCertificate.redundancySeven_count_exhaustion}%
\uses{thm:r6,cor:one-column-extensions}%
\imports{ball-de-beule-mds,duer,seroussi-roth,wu-ding-chen}%
\evidence{certificate-r7}%
\proves{redundancy-seven-classification}%
Certificate R7 classifies the split-free syndromes for
\[
q\in\{7,8,9,11,13,16,17,19,23,25,27,29,31,32\}.
\]
For every prime power $q\geq13$,
the split-free syndromes
at redundancy seven are the persistent locus of size
\(q(q+1)^2/2\), together with the
central nucleus point when $q=2^m$ and $m$ is odd.  At $q=7,8,9,11$ the additional
split-free orbit-size profiles, after removing the persistent locus
and also the central nucleus singleton at \(q=8\), are:
An exponent records multiplicity, so, for example, \(84^5\) means five
distinct orbits of size \(84\).
\[
\begin{array}{c|l}
7&56,\ 84^5,\ 112^2,\ 168^{45},\ 336^{141}\\
8&63,\ 72,\ 84^3,\ 168^4,\ 252^{24},\ 504^{86}\\
9&180^3,\ 240^6,\ 360^{18},\ 720^{27}\\
11&264^2,\ 440,\ 660^2 .
\end{array}
\]
This is the complete split-free classification for all prime powers
\(q\geq7\).  Together with the radius and exact-distance statements below,
it yields the complete code deep-hole classification.
The covering radius is six for \(q=7,9\) and for every \(q\geq11\).
At $q=8$, where
\cite[Thm.~17(1)]{WuDingChen2023} gives radius seven, the diagonal tangent
orbit (nine directions) and the fixed central nucleus are exactly the
distance-seven rows; exhaustive replay puts all 45 other frozen representatives
and the remaining three persistent representatives at distance six. Thus the
$q=8$ classification is complete. The
persistent orbit law is $T/T^6$ modulo inversion and Frobenius, with
tangent splitting in characteristics two and three.
\end{theorem}

\begin{proof}
For fixed irreducible quadratic \(Q\), the persistent fiber has
\(q+1\) points, while for \(Q=\lambda^2\) it has \(q\) non-rank-one
points.  There are \(q(q-1)/2\) irreducible quadratics and \(q+1\)
rational squares.  Hence the persistent count is
\[
 \frac{q(q-1)}2(q+1)+(q+1)q=\frac{q(q+1)^2}{2}.
\]
For \(q\geq37\), choose a first marker outside the at most three
catalecticant intersections, eight collision parameters, and, in
characteristic two, one intersection with the nucleus line.  The first-marker
exclusions put the contracted quartic net outside its persistent intersection;
Lemma~\ref{lem:basepointfree-off-persistent} therefore gives trivial gcd.
Proposition~\ref{prop:r7-pointed} supplies the complete second one-step
package: it excludes the persistent and cyclic/wild carriers before applying
the pointed \(S_3\) cover.
These first-marker bounds and Proposition~\ref{prop:r7-pointed} are
exactly the depth-two data of Theorem~\ref{thm:induction}; the terminal
genus-one curve has at most \(24\) bad rational points.
Corollary~\ref{cor:r7-contained} leaves exactly the persistent locus and
Proposition~\ref{prop:r7-central}.  Certificate R7 covers
\[
q\in\{7,8,9,11,13,16,17,19,23,25,27,29,31,32\}
\]
by the exact primary quotient enumeration.  A separate replay checks the
recorded representatives and their orbits.  A direct-locus replay, independent
of the R7 generator and quotient partition, imports neither: below \(16\) it
constructs the literal four-secant complement, from \(16\) onward it
constructs the classified R6 pointed locus directly, including the
transient marked orbit at \(q=19\), and in every field it intersects all
contraction conditions and rebuilds the complete
\(\PGL_2(q)\)-orbit and Frobenius partition.  It agrees exactly with the
public record in all fourteen fields.

For \(q\geq11\), Seroussi--Roth nonextendability and Dür's equivalence give
\(\rho(\PRS(q-6))=6\).  It remains to close \(q=7,9\).  Here
\(\PRS(q-6)\) has parameters \([q+1,q-6,8]_q\).  If a syndrome \(f\)
had weight seven, adjoining \(f\) to its parity-check matrix would, by
\eqref{eq:extension-distance}, define an
\([q+2,q-5,8]_q\) MDS code.  At \(q=7\) this would be an
\([9,2,8]_7\) MDS code, whose generator columns would give nine distinct
points of \(\PP^1(\F_7)\), although that line has only eight points.  At \(q=9\)
it would be an \([11,4,8]_9\) MDS code, excluded by
Ball--De Beule~\cite[Thm.~5.2]{BallDeBeule2012}: here
\(9=3^2\) and \(4=2\cdot3-2\), the endpoint of their theorem.
Thus every syndrome has weight at most six in both fields.  The nonempty
split-free lists above give the reverse inequality, so the radius is six.
Together with the cited \(q=8\) radius and exact-distance replay, this promotes
the split-free classification over every prime power \(q\geq7\).
\end{proof}

For the four finite rows Certificate R7 contains the canonical
representative, stabilizer, persistent/central flags, and Frobenius link
of every orbit.  The displayed size profile is only a compact summary;
repeated sizes are distinguished by those canonical records.  There
are respectively \(194,119,54,5\) exceptional representatives at
\(q=7,8,9,11\); the complete lists are printed in
\texttt{supplement/CLASSIFICATION-RECORDS.md}, while the matching JSON
is the machine-readable record.  The following gives one printed
canonical representative for every orbit-size block in the displayed
profile; repeated orbits within a block are distinguished in the full
list:
\[
\begin{array}{c|c|l}
q&|\mathcal O|&\text{representative}\\ \hline
7&56 &[0:0:0:0:1:0:0]\\
 &84 &[0:1:0:0:0:1:0]\\
 &112&[0:0:1:0:0:2:0]\\
 &168&[0:0:0:0:1:0:1]\\
 &336&[0:0:0:1:0:1:1]\\
8&63 &[0:0:0:1:0:0:1]\\
 &72 &[0:0:0:0:1:0:0]\\
 &84 &[0:1:1:1:1:3:6]\\
 &168&[1:0:1:1:3:7:7]\\
 &252&[0:0:1:0:1:0:0]\\
 &504&[0:0:1:0:1:0:1]\\
9&180&[0:1:0:0:0:4:0]\\
 &240&[0:0:1:0:1:1:1]\\
 &360&[0:1:0:1:3:3:2]\\
 &720&[0:0:1:1:0:2:2]\\
11&264&[0:1:0:0:0:0:2]\\
  &440&[1:0:1:1:4:3:3]\\
  &660&[1:0:1:0:1:4:8].
\end{array}
\]
For \(q=8,9\), the coordinates use the same polynomial-basis encoding
specified below Table~\ref{tab:r6-exceptional}.

The \(q=19\) lower net
\[
                 W=\langle1,t^3,t^4\rangle
\]
shows why the parameter must remain marked.  Its split squarefree
members are exactly
\[
 U(T^3-uU^3),\qquad
 u\in(\F_{19}^{\times})^3=\{1,7,8,11,12,18\}.
\]
All six contain the marked point at infinity.  The obstruction is
pointed rather than intrinsic to \(W\): at the marker \(0\), for
example, \(U(T^3-U^3)\) is an admissible member.  No coherent sextic
polar line lies in the infinity-pointed orbit.  Thus the R7 proof
escapes over the marker parameter; classifying unpointed bad lower
fibers would retain a false exception.

\begin{table}[ht]
\centering
\small
\caption{Field-range closure ledger.  ``Certificate'' means a complete
finite classification with canonical representatives, not merely an
orbit count.}
\label{tab:field-ranges}
\begin{tabular}{lll}
\toprule
Result & prime powers & closure mechanism\\
\midrule
R5 & $7,8,9,11,13,17,19$ & direct census and independent replay\\
   & $16$ & characteristic-two branch budget,
       Corollary~\ref{cor:r5-binary-shallow};\\
   & & certificate row retained as a regression check\\
   & $q\geq23$ & cubic-cover geometry, Lemma~\ref{lem:s3}\\
R6 & $7,8,9,11,13,16$ & direct syndrome/radius census\\
   & $17,19,23,25,27$ & finite structural bridge certificate\\
   & $q\geq29$ & marked-cover point bound and contained components\\
R7 & $7,8,9,11$ & direct split-free census\\
   & $13,16,17,19,23,25,27,29,31,32$ &
       coherent-polar finite bridge certificate\\
   & $q\geq37$ & marked-cover point bound and contained components\\
\bottomrule
\end{tabular}
\end{table}

The R7 rows classify split-free syndromes.  The covering-radius argument in
Theorem~\ref{thm:r7} promotes them to code deep-hole directions over every
listed field.  The public certificate schema and exact
replay boundaries are given in the supplement; the table does not
promote orbit-size profiles to identifiers.

There is no prime power in any of the open numerical intervals
\((19,23)\), \((27,29)\), or \((32,37)\).  Thus adjacent finite and
geometric rows in Table~\ref{tab:field-ranges} leave no unlisted field
order.
